% SOURCE_AUTHORITY: sga5_fr_workpass.tex, lines 14384--14403 (LF-normalized unit).
% SOURCE_SHA256: 791F4EFFC5E02832D5D77ED03518C8156D6F07E4C8238B03545DB93D883FBB28
Demostremos a). El compuesto $\pi_{X/S}\circ\Fr_{X/S}=\fr_X$ es integral, sobreyectivo y radicial de manera evidente; resulta que $\Fr_{X/S}$ es radicial (EGA I 3.5.6 (ii)); además, $\pi_{X/S}$ es separado y radicial, pues se deduce de $\fr_S$ mediante el cambio de base $X\to S$; resulta que $\Fr_{X/S}$ es integral (EGA II 6.1.5 (v)) y sobreyectivo.

La demostración de b) es evidente, observando que
\[
\JJ\OO_{(S,\fr_S)}=\JJ^p
\]
para todo ideal $\JJ$ de $\OO_S$.

Demostración de c). Se sabe que, para que un morfismo sea un monomorfismo, respectivamente un isomorfismo, es necesario y suficiente que sea radicial y no ramificado, respectivamente radicial, étale y sobreyectivo; de ahí la equivalencia entre (ii) y (iii), respectivamente entre (ii') y (iii'). Queda por probar que (i) equivale a (ii) y (i') a (ii').

Para ello se utiliza el criterio diferencial de no ramificación (SGA I 3) y el isomorfismo
\[
\Omega^1_{X/S}\simeq \Omega^1_{X/X^{(p)}}
\]
dado por la sucesión exacta canónica
\[
\Fr_{X/S}^*(\Omega^1_{X^{(p)}/S})\longrightarrow \Omega^1_{X/S}
\longrightarrow \Omega^1_{X/X^{(p)}}\longrightarrow 0
\]
(SGA II 4), en la cual la primera flecha es visiblemente nula, pues $\Omega^1_{X^{(p)}/S}$ está engendrado por los diferenciales de las secciones de $\OO_{X^{(p)}}$, y el diferencial $d_{X/S}$ se anula sobre $(\OO_X)^p$ y sobre $\OO_S$. Así, (i) es equivalente a (ii).
